\section{An inherited density theorem, independently rechecked}
\label{app:density-audit}

This appendix revalidates a completed result in the supplied archive
\href{https://github.com/VladimirReshetnikov/ProveIt/blob/5cc267af75a635199ac5d93a7d49cbce521daefc/docs/incoming/grh_transfers_density_gaps_and_zero_tests.zip}{\emph{Geometric progressions and nonabelian zero tests}},
specifically its members \texttt{integers.tex} and \texttt{integer\_more.tex}
under \texttt{grh\_transfers/}. It is prior work in this corpus, not a discovery
of the present investigation. The proof is independent of RH and GRH.

For $E\subseteq\mathbb N$, let $d(E)$ be its natural density when it exists,
and let $\overline d(E)$ be its upper density. Call $E$ admissible if it
contains no three distinct terms $a,ar,ar^2$ with $r>1$. Treat separately
the conventions $r\in\mathbb Q$ and $r\in\mathbb Z$. Write
$(b,u)=(\alpha,\overline\alpha)$ in the rational convention and
$(b,u)=(\beta,\overline\beta)$ in the integer convention, where $b$ is the
supremum of existing natural densities and $u$ the supremum of upper densities.

\paragraph{Revalidated theorem.}
In both conventions,
\[
                         u-b>\frac{404}{10^6}.
\]
Thus both strict comparisons posed in McNew's Questions~7.1 and~7.3 hold.
The source report gives stronger constants, recorded below.

\paragraph{Proof.}
Fix an integer $Y\ge12$. Let $S$ be the positive integers whose prime factors
are at most $Y$, including $1$, and set
\[
 c=\prod_{p\le Y}(1-p^{-1}),\qquad
 v=\prod_{p>Y}(1-p^{-2}).
\]
Let $J(t)$ be the maximum size of an admissible subset of $S\cap[1,t]$.
Define two different optimizations:
\[
 B=c\sup_{\substack{A\subseteq S\\A\ {\rm admissible}}}
       \sum_{h\in A}\frac1h,
 \qquad
 U=c\int_1^\infty J(t)\frac{dt}{t^2}.
\]
Both are finite, since
$c\sum_{h\in S}h^{-1}=1$, and hence $U\le1$.

Every positive integer is uniquely $gh$ with $h\in S$ and $g$ having no prime
factor at most $Y$. For a fixed admissible $E$, each fiber
$\{h\in S:gh\in E\}$ is admissible. Consequently
\[
 \sum_{\substack{n\le X\\n\in E}}\frac1n
 \le \frac Bc\sum_{\substack{g\le X\\p\mid g\Rightarrow p>Y}}\frac1g
 =B\log X+O_Y(1).
\]
The last equality follows from periodicity of the rough condition and partial
summation. If $d(E)=d$, the left side is $(d+o(1))\log X$. Thus $b\le B$;
no density assumption on individual fibers is being made.

The squarefree rough integers have density $cv$, by the Chinese remainder
theorem and the summable tail of $\sum_p p^{-2}$. Restrict to these rough parts.
In any geometric progression, the valuations at a prime $p>Y$ form an
arithmetic progression in $\{0,1\}$ and therefore are equal. All three terms
must consequently have the same rough part. This permits independent choices
of admissible smooth fibers.

Choose integers $X_j$ with $X_j/X_{j-1}\to\infty$. For each squarefree rough
$g\in(X_{j-1},X_j]$, permanently assign a maximizing fiber in
$S\cap[1,X_j/g]$. The union of all these fibers, multiplied by their rough
parts, is one admissible set. For fixed $K>1$, at $X=X_j$ its newly assigned
rough parts in $(X/K,X]$ contribute asymptotically
\[
 cvX\int_1^K J(t)\frac{dt}{t^2}.
\]
Indeed, $J$ is a step function with finitely many jumps on $[1,K]$ and the
rough parts have density $cv$. Let first $j$, then $K$, tend to infinity.
This proves $u\ge vU$.

All integers through $12$ belong to $S$. A subset maximizing its cardinality
at every prefix through $12$ would be forced, inductively, to make every
decision of the increasing greedy construction. That greedy prefix is
\[
             \{1,2,3,5,6,7,8,10,11\},
\]
with nine elements. But
\[
             \{1,2,5,6,7,8,9,10,11,12\}
\]
is admissible and has ten. To check this directly, the complete rational-ratio
list through $12$ is
\[
 (1,2,4),\ (1,3,9),\ (2,4,8),\ (3,6,12),\ (4,6,9);
\]
only the last triple has noninteger ratio. Hence every fixed admissible
$A\subseteq S$ loses at least one element against $J(k)$ at some $k\le12$.
That loss persists on $[k,k+1)$, giving
\[
 \int_1^\infty\bigl(J(t)-|A\cap[1,t]|\bigr)\frac{dt}{t^2}
 \ge\frac1{k(k+1)}\ge\frac1{156}.
\]
After taking the supremum over $A$, this yields $U-B\ge c/156$. Therefore
\[
 u-b\ge vU-B
 \ge \frac c{156}-(1-v)
 \ge \frac1{156}\prod_{p\le Y}(1-p^{-1})-\sum_{p>Y}p^{-2}.
\]
At $Y=2000$, the independent exact checker supplied with this report proves
\[
 \frac1{156}\prod_{p\le2000}(1-p^{-1})
 -\sum_{2000<p\le100000}p^{-2}-\frac1{100000}
 >\frac{404}{10^6}.
\]
The omitted prime tail is at most
$\sum_{n>100000}n^{-2}\le1/100000$. This proves the theorem.
\hfill$\square$

\paragraph{Computation and stronger inherited constants.}
The independent checker enumerates all $2^{12}$ subsets and uses the
parametrization $t(b^2,ab,a^2)$, $\gcd(a,b)=1$, for increasing rational
progressions. It confirms the exact minimum prefix-integral loss $1/156$ in
both conventions. Trial division enumerates the primes, and all estimates
use integers and rational numbers.
The two supplied larger checkers were also inspected and rerun successfully.
They establish
\[
 \overline\alpha-\alpha>0.002836,
 \qquad
 \beta\le\frac{5245322990569}{7001316000000}<0.749192.
\]
Together with McNew's published $\overline\beta>0.818410$, the latter gives
$\overline\beta-\beta>0.069218$. The larger computations were replays of the
supplied algorithms; the new independent implementation certifies the simpler
twelve-prefix theorem above.

\paragraph{Literature scope.}
McNew's
\href{https://www.nathanmcnew.com/GPFsets.pdf}{original manuscript},
page~15, states Questions~7.1 and~7.3 in precisely these terms. A targeted
search of current primary sources and his
\href{https://www.nathanmcnew.com/math.html}{research bibliography}
found no later resolution of those two questions. His
\href{https://arxiv.org/html/1808.04923v2}{2021 divisor-graph paper},
Theorem~9, concerns optimal finite subsets; the
\href{https://www.nathanmcnew.com/w45.pdf}{2022 Cao--Fang--McNew paper}
concerns greedy interval endpoints. Neither is the compatibility statement
proved here. This search does not certify worldwide priority. Rankin's
natural-density optimality conjecture and the comparison $\beta>\alpha$
remain unresolved by this argument. The connection to \emph{grh-lean} is the
earlier transfer program's use of additive restrictions on prime exponents;
the present proof additionally uses classical smooth--rough decomposition.
